
\documentclass[12pt]{article}

\usepackage[hyperfootnotes=false]{hyperref}

\makeatletter
\renewcommand\@makefnmark{%
  \hbox{\textsuperscript{\textcolor{blue}{\@thefnmark}}}%
}
\renewcommand\@makefntext[1]{%
  \parindent 1em\noindent
  \hbox{\textsuperscript{\@thefnmark}} #1%
}
\makeatother

\usepackage[margin=0.85in,footskip=0.25in]{geometry}
\newcommand*\dif{\mathop{}\!\textnormal{\slshape d}}


\usepackage{tcolorbox}
\usepackage{physics}
\usepackage{amsfonts}
\usepackage{lipsum}
\usepackage{amsmath,amssymb,amsthm}
\usepackage{mathtools}
\usepackage[dvipsnames]{xcolor}
\usepackage{multirow}

\usepackage{xcolor} 
\definecolor{grey}{rgb}{0.9,0.9,0.9}
\renewcommand{\baselinestretch}{1.2} 

\usepackage[scaled]{helvet}
\usepackage[T1]{fontenc}

\usepackage{subcaption}
\usepackage{graphicx, float}
\usepackage{wrapfig}
\usepackage{caption}
\usepackage{hyperref}
\usepackage{relsize}

\usepackage{asymptote}
\setlength{\parindent}{0pt}

\usepackage{mdframed}
\newmdenv[
  linecolor=black,
  linewidth=1pt,
  topline=false,
  bottomline=false,
  rightline=false,
  innertopmargin=0.5\baselineskip,
  innerbottommargin=0.5\baselineskip,
  innerleftmargin=10pt
]{problem}

\newcommand{\dxdt}{\frac{dx}{dt} }
\newcommand{\half}{\frac{1}{2}}
\newcommand{\cell}[2][1.25cm]{\makebox[#1]{\ensuremath{#2}}}
\newcommand{\scell}[2][1cm]{\makebox[#1]{\ensuremath{#2}}}
\newcommand{\rtwo}{\mathbb{R}^2}
\newcommand{\rthree}{\mathbb{R}^3}


\newcommand{\vect}[1]{\begingroup\renewcommand{\arraystretch}{0.5}\mathbf{#1}\endgroup}


\usepackage{xstring}
\newif\ifhasneg

\newcommand{\realtwovector}[2]{%
  \begingroup
  \renewcommand{\arraystretch}{0.75}%
  \hasnegfalse
  \IfBeginWith{#1}{-}{\hasnegtrue}{}%
  \IfBeginWith{#2}{-}{\hasnegtrue}{}%
  \ifhasneg
    \begin{bmatrix}
      \IfBeginWith{#1}{-}{#1}{\phantom{-}#1} \\[0pt]
      \IfBeginWith{#2}{-}{#2}{\phantom{-}#2}
    \end{bmatrix}%
  \else
    \begin{bmatrix}#1 \\ #2 \end{bmatrix}%
  \fi
  \endgroup
}

\newcommand{\realthreevector}[3]{%
  \begingroup
  \renewcommand{\arraystretch}{0.4}%
  \begin{bmatrix*}[r]
    \scriptstyle #1 \\[2pt]
    \scriptstyle #2 \\[2pt]
    \scriptstyle #3
  \end{bmatrix*}%
  \endgroup
}
\newcommand{\vmag}[1]{\left\lVert \vect{#1} \right\rVert}
\newcommand{\proj}[2]{\text{proj}_{\vect{#1}} \vect{#2}}

\date{March 27, 2026}
% \date{November 10, 2025\\ \small Revised: November 24, 2025}
\title{%
  LA Assessment - Least Squares Approximation 
}
\author{Hugo Hu}

\begin{document}
\maketitle

\section*{Problem}

Yellow cab fares in NYC are usually metered, not flat-fare. Could we use historical trip records to build a regression
model to predict fares for future rides?

\section*{Data and Technical Information}
The NYC Taxi and Limousine Commission provides trip record data recording information about every taxi trip made in the city. These datasets are available for each month, located 
at\\ \url{https://www.nyc.gov/site/tlc/about/tlc-trip-record-data.page}. \\

As a training set, I used twelve months of data (December 2024 to November 2025), randomly selecting half of all entries for inclusion. Of these trips, 
fixed-fare and other special rate trips were excluded, leaving a total of over 17 million rows. \\

The regression model uses information like the time of day, day of week, the month, and the start and end zones, as a part of its "known data" matrix. 
But not all data available is used, because we only wish to consider pieces of information that can be reasonably identified before beginning a trip. For example,
a rider would not know the total distance driven, nor total time in transit, before starting. 
\\

\newpage
\section*{One-Hot Encoding and Dependency}
In a linear regression, we work with numbers. Unfortunately, not everything can be easily encoded as a number! For example, consider the challenge of encoding
a month as a number: it seems simple to designate January through December as $1$ through $12$.\\

But this isn't really the case! Months do not exist on a line (more so a circle, it could be argued). The fix for this is "one-hot encoding". Instead of creating
a single column to populate with an encoded month value, we create $12$ columns each representing a month. Almost like a boolean datatype, we use $1$ or $0$ to 
represent whether it is that month. \\

For example, to represent December in this encoding, we'd have: $\left\{0, 0, 0, 0, 0, 0, 0, 0, 0, 0, 0, 1\right\}$. \\

But knowing that we'll find inverses later, it seems like this technique has a flaw that creates possible dependency. In the encoding above, isn't it true that we
only need $11$ pieces of information to represent all $12$ states?\\

If we knew that it wasn't January through November, we'd conclude that it must then be December. 
To avoid this dependency, otherwise known as a "dummy variable", we will deliberately drop a column from each feature (month, time of day, day of week, etc). \\

\section*{Linear Algebra Setup}
Our problem is simple: solve for vector $\vect{w}$ in $X \vect{w} = \vect{y}$. \\

$X$ is the data matrix, and for each ride, includes the following information:
\begin{itemize}
    \item Pickup Zone (262)
    \item Drop-off Zone (262)
    \item Pickup Hour (24)
    \item Day of Week (7)
    \item Month (12)
\end{itemize}

\newpage

After hot-encoding and dropping a column, our matrix looks roughly like: \\
\begin{tabular}{|p{2.5cm}||p{2.5cm}|p{2.5cm}|p{2.5cm}|p{2.5cm}|p{2.5cm}|}
\hline
\multicolumn{6}{|c|}{Matrix $X$ of size 17M x 561} \\
\hline
Ride Data \# & Pickup Zone (261) & Drop-off Zone (261) & Pickup Hour (23) & Day of Week (6) & Pickup Month (11) \\
\hline
Ride 1 & one-hot & one-hot & one-hot & one-hot & one-hot \\
\hline
Ride 2 & one-hot & one-hot & one-hot & one-hot & one-hot \\
\hline
... & ... & ... & ... & ... & ... \\
\hline
Ride 17M & one-hot & one-hot & one-hot & one-hot & one-hot \\
\hline
\end{tabular} \newline\newline



Our vector $\vect{y}$ is a vertical $17$M x $1$ vector, where each element corresponds to the pre-tip fare of that associated ride. This is the data we'd like 
to be able to infer.\\

Lastly, $\vect{w}$ is the coefficient vector of size $561$ x $1$. We are looking to find $\vect{w}$. \\

\section*{Workthrough}
Though these computations were performed in numpy, the following work is equivalent to what was implemented in code. \\

We begin with:\\
$X\vect{w} = \vect{y}$ \\
We know that $X$ is not a square matrix, and thus, we cannot find its matrix. Let's change that!\\

$X^T X \vect{w} = X^T \vect{y}$ \\
For any matrix $M$, $M^TM$ is a square matrix (with an inverse), so $X^TX$ has an inverse.\\

$(X^T X)^{-1} (X^T X) \vect{w} = (X^T X)^{-1} (X^T) \vect{y}$\\

$\vect{w} = (X^TX)^{-1} (X^T) \vect{y}$. 

\newpage
\section*()
\end{document}
